\section{Causal Reasoning: Why a Correlation Is Not a Strategy}

\subsection{The Anatomy of False Strategic Intuitions}

Founders are bombarded with industry maxims that masquerade as causal truths when they are merely superficial correlations:
\begin{itemize}[leftmargin=10pt,itemsep=1pt,topsep=1pt]
  \item \textit{Elite schools hire celebrity professors and charge ₹10 lakh in tuition; therefore, hiring celebrity professors will allow us to charge ₹10 lakh.}
  \item \textit{Billion-dollar enterprise software companies deploy 50-person outbound sales armies; therefore, hiring an outbound sales army will generate enterprise hypergrowth.}
  \item \textit{Prominent venture-backed founders work 90-hour weeks; therefore, working 90-hour weeks guarantees market triumph.}
\end{itemize}

Every one of these inferences commits elementary causal errors: selection bias, reverse causality, and omitted variable bias. Institutional prestige causes both celebrity hiring and pricing power; surging enterprise revenue permits the hiring of aggressive sales teams rather than resulting from it; and work stamina is a prerequisite that cannot compensate for a product nobody wants.

The founder's central operational question is never \textit{what correlates with past success?} The only question that matters for capital allocation is: \textbf{What precise outcome will change if I intervene directly in this system?}

\subsection{Mapping the Full Causal Graph}

Consider a simple linear hypothesis:
\begin{center}
\small
\begin{tabular}{c}
\textbf{Experienced Faculty} $\;\longrightarrow\;$ \textbf{Superior Pedagogy} \\
$\downarrow$ \\
\textbf{Student Mastery} $\;\longrightarrow\;$ \textbf{Referral Growth}
\end{tabular}
\end{center}

This naive graph ignores the real underlying dynamics:
\begin{itemize}[leftmargin=10pt,itemsep=1pt,topsep=1pt]
  \item Highly motivated, elite students proactively seek out prestigious faculty (selection).
  \item An established institutional brand attracts both top faculty and top corporate recruiters (omitted variable).
  \item Corporate hiring managers frequently filter candidates by university pedigree rather than verified technical competence (signaling distortion).
\end{itemize}

If a new school spends its entire treasury hiring a famous academic while admitting students who lack foundational math, the chain snaps immediately. The pedagogy fails, student outcomes collapse, and referral growth never materializes.

\subsection{The 3 Causal Interrogation Questions}

Before deploying capital or changing strategic direction, run your proposal through 3 rigorous inquiries:

\begin{enumerate}[leftmargin=12pt,itemsep=2pt,topsep=1pt,parsep=0pt]
  \item \textbf{The Intervention Question:} \textit{If we alter this specific operational variable while holding all surrounding conditions stable, what exact metric should shift, by how much, and on what precise timeline?}
  \item \textbf{The Rival Explanation Question:} \textit{What alternative operational mechanisms, selection biases, or macroeconomic tailwinds could produce the exact same observable outcome?}
  \item \textbf{The Mechanism Question:} \textit{What is the complete physical and psychological chain connecting our capital expenditure to customer value creation, and at which specific node is that chain most likely to break?}
\end{enumerate}

\subsection{Triangulation Over Naive Empiricism}

Early-stage ventures operate in environments too noisy and sample-constrained for clean randomized controlled trials. Running an A/B test with 30 website visitors produces statistical noise, not scientific certainty.

The defense against self-delusion is triangulation: combining bounded pilot interventions, rigorous customer discovery interviews, unit-level operational telemetry, and mechanism-level causal logic.

\textbf{A founder must become allergic to any strategic recommendation whose validity relies exclusively on observing the celebratory habits of visible winners.}
